\documentclass[12pt]{extarticle}
\def\activityid{F02-FAC-v2}\usepackage[letterpaper,margin=.65in,top=.60in,bottom=.65in]{geometry}
\usepackage[T1]{fontenc}\usepackage[default]{sourcesanspro}
\usepackage{microtype,tikz,array,tabularx,fancyhdr,amsmath,amssymb}
\usepackage[hidelinks]{hyperref}\usetikzlibrary{calc,arrows.meta}
\setlength{\parindent}{0pt}\setlength{\parskip}{7pt}\setlength{\headheight}{0pt}\setlength{\footskip}{26pt}
\setlength{\emergencystretch}{2em}\pagestyle{fancy}\fancyhf{}\renewcommand{\headrulewidth}{0pt}\renewcommand{\footrulewidth}{.4pt}
\fancyfoot[L]{\scriptsize Bellingham Math Circle / Week 2 / \activityid}\fancyfoot[R]{\scriptsize \thepage}
\definecolor{ink}{gray}{.12}\definecolor{soft}{gray}{.95}\color{ink}
\newcommand{\PageTitle}[2]{{\fontsize{10}{12}\selectfont\bfseries WEEK 2\quad TOGGLE LAB\hfill #1\par}\vspace{3pt}{\fontsize{26}{29}\selectfont\bfseries #2\par}\vspace{2pt}}
\newcommand{\NameLine}{{\small Name: \rule{2.65in}{.4pt}\hfill Date: \rule{1.35in}{.4pt}\par}\vspace{4pt}}
\newcommand{\Task}[2]{\par\vspace{3pt}{\large\bfseries #1}\par #2\par}
\newcommand{\Subhead}[1]{\par\vspace{3pt}{\large\bfseries #1}\par}
\newcommand{\RuleBox}[1]{\par\begingroup\setlength{\fboxsep}{8pt}\colorbox{soft}{\parbox{\dimexpr\linewidth-16pt\relax}{#1}}\endgroup\par}
\newcommand{\WriteBox}[2]{\par\textbf{#1}\par\vspace{2pt}\begin{tikzpicture}\draw[black!40,line width=.5pt] (0,0) rectangle (\dimexpr\linewidth-.5pt\relax,#2);\end{tikzpicture}\par}
\newcommand{\StudentFooter}{\vfill{\small Try it with objects. Explain by pointing, talking, or drawing.}\par}
\newcommand{\LampRing}[3]{\begin{tikzpicture}[x=1in,y=1in]
\foreach \i in {1,...,#1}{\coordinate (v\i) at ({90-360*(\i-1)/#1}:#2);}
\foreach \i in {1,...,#1}{\pgfmathtruncatemacro{\j}{mod(\i,#1)+1}\draw[thick] (v\i)--node[midway,fill=white,inner sep=2pt,font=\small]{\i\j}(v\j);}
\foreach \i in {1,...,#1}{\node[circle,draw,fill=white,minimum size=.52in,inner sep=0pt,font=\bfseries] at(v\i){\i};}
\foreach \i in {#3}{\node[circle,draw,fill=ink,text=white,minimum size=.52in,inner sep=0pt,font=\bfseries] at(v\i){\i};}
\end{tikzpicture}}
\newcommand{\Lights}[1]{\begin{tikzpicture}[x=.55in,y=.55in,baseline=-.10in]\foreach \v [count=\i] in {#1}{\ifnum\v=1\fill (\i,0) circle(.16in);\else\draw (\i,0) circle(.16in);\fi}\end{tikzpicture}}
\newcommand{\ClockRing}[2]{\begin{tikzpicture}[x=1in,y=1in]\draw[black!30] (0,0) circle(#2);\pgfmathtruncatemacro{\n}{#1-1}\foreach\i in {0,...,\n}{\node[circle,draw,fill=white,minimum size=.35in,inner sep=0pt] at({90-360*\i/#1}:#2){\i};}\draw[-{Stealth},thick] (-.35,.15) arc(155:25:.38);\node[font=\small] at(0,-.18){clockwise};\end{tikzpicture}}
\newcommand{\Key}[2]{\begin{center}\renewcommand{\arraystretch}{1.55}\begin{tabular}{c|*{#1}{c}}in&#2\end{tabular}\end{center}}

\begin{document}

\PageTitle{FACILITATOR / 1}{Prepare the toggle lab}
\RuleBox{The mathematical question is reachability, then optimization: which lamp patterns can local pair-flips produce, and how cheaply? This replaces the old missing-stroke core. Do not mix the old union rule with toggling.}
\Subhead{Materials / seven children / two adults}
Bring \textbf{29 two-sided counters}: four per K--1/middle child (24), five for upper; 8 spares; pencils and scratch paper. Paper circles with one shaded face work. Tap edges with fingers; no electronics. Print 3 K--1, 3 middle, 1 upper, 1 extra reserve, and this guide. Lend a spare for the six-ring.
\Subhead{Readiness and entry points}
\textbf{K--1:} no independent reading; match two counter faces; count or pair up to four. Adult reads and models both endpoints flipping, including ON to OFF. Reasoning: undoing and an invariant.\par
\textbf{2--3:} labels 1--4, counting to four, odd/even or visible pairing; read-aloud works. Reasoning: classify cases, organize all possibilities, explain interior cancellation.\par
\textbf{4--5:} count to six, track yes/no choices, compare lengths. Reasoning: necessity versus sufficiency, exhaustive two-branch argument, construction plus lower bound.\par
\textbf{Extra:} read a graph, parity, follow paths; no algebra notation required. The abstraction is one theorem covering all finite graphs, then a unique solution on trees.
\Subhead{A shared 60-minute hour}
\textbf{0--10:} explore counters and make pictures freely. \textbf{10--15:} demonstrate one edge press and its undo. Launch: ``Can we make any picture we want?'' \textbf{15--35:} K pair targets; middle target collection; upper cheapest solutions. Adult A anchors K; organizer visits middle and upper. \textbf{35--40:} movement/reset: stand, stretch, leave work ready to return to. \textbf{40--55:} continue the collection or use the optional proof paths below; offer the six-ring/extra only if wanted. \textbf{55--60:} share one discovery or impossibility; tidy counters.
Give one page at a time. Rotate mover/target-maker/checker in triplets, each retaining a mat and counters. Schedule an adult visit to hear the upper child's optimality argument.
\Subhead{Hints that leave room to discover}
Check that ON flips OFF. Ask what a repeated press does. For parity, compare both OFF/both ON/mixed endpoints; for paths, count touches at a middle vertex. Cover later table columns and decide one edge at a time. A stalled child can stay on four lamps.
\Subhead{Satisfying stops / optional proof}
\textbf{K--1 stop:} make and undo two pair targets. \textbf{Optional:} explain the impossible singleton by pairing ON lamps.\par
\textbf{2--3 stop:} collect eight four-lamp targets. \textbf{Optional:} prove odd targets fail and paths make every pair.\par
\textbf{4--5 stop:} solve both five-ring targets minimally; explain why one press fails. \textbf{Optional:} show each first-edge choice forces the rest, then use complements to prove the two-press bound.
\newpage
\PageTitle{FACILITATOR / 2}{Parity and all the pictures}
\Subhead{K--1 solutions and explanations}
Two identical presses undo each other; replaying any list in reverse undoes it. Top/right (1,2) needs 12. Top/bottom (1,3) uses 12 then 23. A single ON lamp is impossible under the pair rule. Both OFF endpoints become ON (+2), both ON become OFF (-2), or one ON trades places (0). Starting at zero, the number ON stays even. Children may demonstrate these three cases and pair the counters instead of saying ``parity.'' The new one-lamp button breaks this invariant and makes the singleton immediately.
\Subhead{Middle: the complete set}
The eight states, listed by their ON lamps, are:
\begin{center}\begin{tabular}{c|l}Number ON&Possible sets\\\hline 0&none\\2&12, 13, 14, 23, 24, 34\\4&1234\end{tabular}\end{center}
Here ``13'' in the \emph{state list} means lamps 1 and 3, not a legal edge. Avoid conflating a target list with a move list. All six pairs can be made by a path on the square. Each interior vertex is touched twice and returns OFF; each end is touched once. To make all four, use edges 12 and 34. No odd-size state is possible. This proves both that the list is complete and that every listed state is reachable. There are 12 blank recording boxes so their count does not give away the answer.
\Subhead{Why paths generalize}
Pair the desired ON vertices. For each pair, choose a path and press its edges. At an intermediate vertex, contributions cancel in pairs; endpoints give one flip. Shared edges or vertices cause no problem: all flip counts add modulo two. An edge pressed twice can be deleted from the list. On any connected graph, exactly the even targets are reachable. Connectivity matters; see the extra.
\Subhead{What is being added?}
Old shape addition used union: a shared stroke was kept once. Toggling is \textbf{symmetric difference}: something present twice cancels. This new operation restores information in the identity \((A\mathbin{\triangle}B)\mathbin{\triangle}B=A\). A one-minute stick-picture comparison is an optional bridge for returners, not a second core activity. Lowell Handout 11, problem 11.3, asked for splitting and recombining shapes; Handout 9, problem 9.1, used door toggles indexed by divisors. Here the adjacent-pair move and reachability problem are new.
\newpage
\PageTitle{FACILITATOR / 3}{Two answers, then an optimum}
\Subhead{Upper page 1: constructions plus lower bounds}
Target \{1,3\}: edges \{12,23\} or \{51,45,34\}. Minimum \textbf{2}. One press can light only adjacent lamps; these are not adjacent. Target \{1,2,3,4\}: \{12,34\} or \{23,45,51\}. Minimum \textbf{2}, because one press changes at most two lamps. Children should supply both the successful move list and the reason one press cannot suffice.
\Subhead{Upper page 2: exactly two reduced solutions}
For target \{1,3\}, the decision table is:
\begin{center}\renewcommand{\arraystretch}{1.4}\begin{tabular}{c|ccccc|c}&51&12&23&34&45&lamp 5\\\hline First no&no&yes&yes&no&no&OFF\\First yes&yes&no&no&yes&yes&OFF\end{tabular}\end{center}
Once 51 is chosen, lamp 1 forces 12, lamp 2 forces 23, and so on. There can be at most two solutions using each edge zero or one times. Complementing an edge set preserves the target, because each lamp has exactly two incident edges: a flip count of 0, 1, 2 becomes 2, 1, 0, with the same parity. Thus every reachable target has exactly two reduced solutions. There are infinitely many unreduced move lists if repeated presses are permitted; state the distinction.
For target \{1\}, first-no forces only 12,23,34,45, leaving lamp 5 ON incorrectly; first-yes forces only 51, again leaving lamp 5 ON. Both branches fail. Parity explains the same obstruction globally.
\Subhead{Upper page 3: the five-ring speed limit}
An even target is reachable by paths. Its two reduced press sets have sizes \(k\) and \(5-k\); at least one is at most 2. Any repeated presses can be cancelled without changing the target, so the smaller reduced set is globally shortest. On an odd ring the two sizes never tie. On the six-ring, opposite lamps 1 and 4 have solutions \{12,23,34\} and \{61,56,45\}, each size 3. The same forced-choice proof shows there is no third reduced set, hence no shorter route. More generally the two cycle solutions have sizes \(k,n-k\).
\Subhead{Undergraduate interpretation}
Write OFF=0 and ON=1. Each edge is a column with two 1's in a vertex-edge incidence matrix. A move set solves \(Ax=t\) over the two-element field. The all-edge cycle is a nonzero null vector; on a single cycle it is the only one. Our forced-choice proof is binary elimination without matrix notation. The connected-graph image is the even-parity subspace; its dimension is \(n-1\).
\newpage
\PageTitle{FACILITATOR / 4}{Networks and real optimization}
\Subhead{Extra: every component matters}
Target \{A,D\} is impossible: the triangle and the separate edge each contain an odd number of target vertices, while every move changes two vertices in one component. Target \{A,C,D,E\} uses AC and DE, minimally two presses. In general, a target is reachable \emph{if and only if each component contains an even number of targets}. Necessity is component parity; sufficiency pairs targets within each component and uses paths. Isolated vertices are components too.
\Subhead{Extra: why a tree has one solution}
Delete an edge e. Count desired ON vertices on either side. Presses wholly on one side affect its parity by zero; the only move crossing the cut is e. Thus e is pressed an odd number of times exactly when that side has an odd target count; otherwise it is pressed an even number of times (possibly zero). The two sides agree because the whole target is even. These decisions uniquely specify every edge. Existence follows from the path construction. They therefore give the unique reduced solution and a shortest move list; repeated edges can only lengthen it. A leaf-removal algorithm gives another constructive explanation.
\Subhead{Research connection and its boundary}
The pressed-edge set has odd degree exactly at the target vertices T: this is a \textbf{T-join}. The children's optimization is exactly minimum-cardinality T-join on a cycle. General networks can have many cycle choices; edge costs lead to the minimum-weight T-join problem and route-inspection algorithms. We prove the elementary parity and tree/cycle cases here, not the general matching algorithm or its polyhedral theorem.
\textbf{Edmonds and Johnson (1973)}, \emph{Matching, Euler tours and the Chinese postman}, \emph{Mathematical Programming} 5, 88--124. The publisher abstract identifies matching, Euler tours, and the optimization result: \href{https://doi.org/10.1007/BF01580113}{doi:10.1007/BF01580113}.\par
\textbf{Cornu\'ejols}, \emph{Combinatorial Optimization: Packing and Covering}, Chapter 2, pp. 27--29, defines T-joins, explains shortest-path/matching reduction, and states the Edmonds--Johnson theorem: \href{https://www.sfu.ca/~mdevos/notes/pack-cover/Cornuejols.pdf}{author's monograph}. It uses acyclic representatives; a cycle can be deleted without changing odd-degree vertices and without increasing nonnegative cost.\par
\textbf{Undergraduate anchor:} MIT 18.06SC, \href{https://ocw.mit.edu/courses/18-06sc-linear-algebra-fall-2011/pages/ax-b-and-the-four-subspaces/graphs-networks-incidence-matrices/}{Graphs, Networks, Incidence Matrices}. Its usual real-valued matrix framework motivates the representation; the binary adaptation and proof above are ours.
\Subhead{Teaching sources and use history}
\emph{Math Circle by the Bay}, preface pp. viii--x, advocates themes with deep mathematical context, manipulatives, varied pace, and reserve challenges. Rozhkovskaya, Lesson 3 ``At the lesson,'' item 1, lets children attempt and explain before a table is introduced; Lesson 7, item 1, warns how missed moves can distort a game investigation. Our staffing and timing are adaptations. Record actual tasks as F02-K/M/U/X-v2; preparation is not evidence of teaching. Reserve the general graph version for a later year if unused.
\textbf{Verification:} \texttt{verify.py} enumerates move subsets for 4-, 5-, and 6-cycles, checks two complementary solutions per even target, minima, and the extra examples. The arguments above stand independently of computation.

\end{document}
